\Chapter{27}

\Verse{1}
{
    \zA{话说黛玉正自悲泣，忽听院门响处，只见宝钗出来了，宝玉袭人一群人都送出 来。 待要上去问着宝玉，又恐当着众人问羞了宝玉不便，因而闪过一旁，让宝钗去
        了，宝玉等进去关了门，方转过来，尚望着门洒了几点泪。 自觉无味，转身回来， 无精打彩的卸了残妆。 紫鹃雪雁素日知道黛玉的情性：无事闷坐，不是愁眉，便是
        长叹，且好端端的不知为着什么，常常的便自泪不干的。 先时还有人解劝，或怕他 思父母，想家乡，受委屈，用话来宽慰。}
}
{
    \eA{Lin Tai-yü, we must explain in taking up the thread of our narrative, was
        disconsolately bathed in tears, when her ear was suddenly attracted by the
        creak of the court gate, and her eyes by the appearance of Pao-ch'ai beyond
        the threshold. Pao-yü, Hsi Jen and a whole posse of inmates then walked out.
        She felt inclined to go up to Pao-yü and ask him a question;}
}
{
}

\Verse{2}
{
    \zA{谁知后来一年一月的，竟是常常如此，把 这个样儿看惯了，也都不理论了。 所以也没人去理他，由他闷坐，只管外间自便去 了。
        那黛玉倚着床栏杆，两手抱着膝，眼睛含着泪，好似木雕泥塑的一般，直坐到 二更多天方才睡了。 一宿无话。 至次日乃是四月二十六日，原来这日未时交芒种节。
        尚古风俗：凡交芒种节的 这日，都要设摆各色礼物，祭饯花神，言芒种一过，便是夏日了，众花皆卸，花神 退位，须要饯行。
        闺中更兴这件风俗，所以大观园中之人都早起来了。}
}
{
    \eA{but dreading that if she made any inquiries in the presence of such a
        company, Pao-yü would be put to the blush and placed in an awkward position,
        she slipped aside and allowed Pao-ch'ai to prosecute her way. And it was
        only after Pao-yü and the rest of the party had entered and closed the gate
        behind them that she at last issued from her retreat.}
}
{
}

\Verse{3}
{
    \zA{那些女孩子 们，或用花瓣柳枝编成轿马的，或用绫锦纱罗叠成干旄旌幢的，都用彩线系了，每 一棵树头每一枝花上，都系了这些物事。
        满园里绣带飘摇，花枝招展，更兼这些人 打扮的桃羞杏让，燕妒莺惭，一时也道不尽。 且说宝钗、迎春、探春、惜春、李纨、凤姐等并大姐儿、香菱与众丫鬟们，都
        在园里玩耍，独不见黛玉，迎春因说道：“林妹妹怎么不见？ 好个懒丫头，这会子 难道还睡觉不成？” 宝钗道：“你们等着，等我去闹了他来。”}
}
{
    \eA{Then fixing her gaze steadfastly on the gateway, she dropped a few tears.
        But inwardly conscious of their utter futility she retraced her footsteps
        and wended her way back into her apartment. And with heavy heart and
        despondent spirits, she divested herself of the remainder of her
        habiliments.}
}
{
}

\Verse{4}
{
    \zA{说着，便撂下众人， 一直往潇湘馆来。 正走着，只见文官等十二个女孩子也来了，上来问了好，说了一 回闲话儿，才走开。
        宝钗回身指道：“他们都在那里呢，你们找他们去，我找林姑 娘去就来。” 说着，逶迤往潇湘馆来。 忽然抬头见宝玉进去了，宝钗便站住，低头
        想了一想：“宝玉和黛玉是从小儿一处长大的，他兄妹间多有不避嫌疑之处，嘲笑 不忌，喜怒无常； 况且黛玉素多猜忌，好弄小性儿，此刻自己也跟进去，一则宝玉
        不便，二则黛玉嫌疑，倒是回来的妙。”}
}
{
    \eA{Tzu Chüan and Hsüeh Yen were well aware, from the experience they had reaped
        in past days, that Lin Tai-yü was, in the absence of anything to occupy her
        mind, prone to sit and mope, and that if she did not frown her eyebrows, she
        anyway heaved deep sighs; but they were quite at a loss to divine why she
        was, with no rhyme or reason, ever so ready to indulge, to herself, in
        inexhaustible gushes of tears.}
}
{
}

\Verse{5}
{
    \zA{想毕，抽身回来，刚要寻别的姊妹去。 忽见面前一双玉色蝴蝶，大如团扇，一 上一下，迎风翩跹，十分有趣。 宝钗意欲扑了来玩耍，遂向袖中取出扇子来，向草
        地下来扑。 只见那一双蝴蝶忽起忽落，来来往往，将欲过河去了。 引的宝钗蹑手蹑 脚的，一直跟到池边滴翠亭上，香汗淋漓，娇喘细细。
        宝钗也无心扑了，刚欲回来， 只听那亭里边嘁嘁喳喳有人说话。 原来这亭子四面俱是游廊曲栏，盖在池中水上， 四面雕镂？ 子，糊着纸。}
}
{
    \eA{At first, there were such as still endeavoured to afford her solace; or who,
        suspecting lest she brooded over the memory of her father and mother, felt
        home-sick, or aggrieved, through some offence given her, tried by every
        persuasion to console and cheer her;}
}
{
}

\Verse{6}
{
    \zA{宝钗在亭外听见说话，便煞住脚往里细听。 只听说道：“你 瞧这绢子果然是你丢的那一块，你就拿着； 要不是，就还芸二爷去。” 又有一个说：
        “可不是我那块!拿来给我罢。” 又听道：“你拿什么谢我呢？ 难道白找了来不成？” 又答道：“我已经许了谢你，自然是不哄你的。”
        又听说道：“我找了来给你，自 然谢我； 但只是那拣的人，你就不谢他么？” 那一个又说道：“你别胡说。 他是个 爷们家，拣了我们的东西，自然该还的。
        叫我拿什么谢他呢？”}
}
{
    \eA{but, as contrary to all expectations, she subsequently persisted time and
        again in this dull mood, through each succeeding month and year, people got
        accustomed to her eccentricities and did not extend to her the least
        sympathy. Hence it was that no one (on this occasion) troubled her mind
        about her, but letting her sit and sulk to her heart's content, they one and
        all turned in and went to sleep.}
}
{
}

\Verse{7}
{
    \zA{又听说道：“你不 谢他，我怎么回他呢？ 况且他再三再四的和我说了，若没谢的，不许我给你呢。” 半晌，又听说道：“也罢，拿我这个给他，算谢他的罢。
        你要告诉别人呢？ 须得起 个誓。” 又听说道：“我要告诉人，嘴上就长一个疔，日后不得好死！” 又听说道：
        “嗳哟!咱们只顾说，看仔细有人来悄悄的在外头听见。 不如把这？ 子都推开了， 就是人见咱们在这里，他们只当我们说玩话儿呢。 走到跟前，咱们也看的见，就别
        说了。”}
}
{
    \eA{Lin Tai-yü leaned against the railing of the bed, clasping her knees with
        both hands, her eyes suffused with tears. She looked, in very truth, like a
        carved wooden image or one fashioned of mud. There she sat straight up to
        the second watch, even later, when she eventually fell asleep. The whole
        night nothing remarkable transpired. The morrow was the 26th day of the
        fourth moon.}
}
{
}

\Verse{8}
{
    \zA{宝钗外面听见这话，心中吃惊，想道：“怪道从古至今那些奸淫狗盗的人，心 机都不错，这一开了，见我在这里，他们岂不臊了？ 况且说话的语音，大似宝玉房
        里的小红。 他素昔眼空心大，是个头等刁钻古怪的丫头，今儿我听了他的短儿，‘人 急造反，狗急跳墙’，不但生事，而且我还没趣。 如今便赶着躲了料也躲不及，少
        不得要使个‘金蝉脱壳’的法子。” 犹未想完，只听“咯吱”一声，宝钗便故意放 重了脚步，笑着叫道：“颦儿，我看你往那里藏！”}
}
{
    \eA{Indeed on this day, at one p.m., commenced the season of the 'Sprouting
        seeds,' and, according to an old custom, on the day on which this feast of
        'Sprouting seeds' fell, every one had to lay all kinds of offerings and
        sacrificial viands on the altar of the god of flowers.}
}
{
}

\Verse{9}
{
    \zA{一面说一面故意往前赶。 那亭 内的小红坠儿刚一推窗，只听宝钗如此说着往前赶，两个人都唬怔了。 宝钗反向他 二人笑道：“你们把林姑娘藏在那里了？”
        坠儿道：“何曾见林姑娘了？” 宝钗道： “我才在河那边看着林姑娘在这里蹲着弄水儿呢。 我要悄悄的唬他一跳，还没有走
        到跟前，他倒看见我了，朝东一绕，就不见了。 别是藏在里头了？” 一面说，一面 故意进去，寻了一寻，抽身就走，口内说道：“一定又钻在山子洞里去了。
        遇见蛇， 咬一口也罢了！”}
}
{
    \eA{Soon after the expiry of this season of 'Sprouting seeds' follows
        summertide, and us plants in general then wither and the god of flowers
        resigns his throne, it is compulsory to feast him at some entertainment,
        previous to his departure. In the ladies' apartments this custom was
        observed with still more rigour;}
}
{
}

\Verse{10}
{
    \zA{一面说，一面走，心中又好笑：“这件事算遮过去了。 不知他二 人怎么样？” 谁知小红听了宝钗的话，便信以为真，让宝钗去远，便拉坠儿道：“了不得了!
        林姑娘蹲在这里，一定听了话去了！” 坠儿听了，也半日不言语。 小红又道：“这 可怎么样呢？” 坠儿道：“听见了，管谁筋疼!各人干各人的就完了。”
        小红道： “要是宝姑娘听见还罢了。 那林姑娘嘴里又爱克薄人，心里又细，他一听见了，倘 或走露了，怎么样呢？”}
}
{
    \eA{and, for this reason, the various inmates Of the park of Broad Vista had,
        without a single exception, got up at an early hour. The young people either
        twisted flowers and willow twigs in such a way as to represent chairs and
        horses, or made tufted banners with damask, brocaded gauze and silk, and
        bound them with variegated threads. These articles of decoration were alike
        attached on every tree and plant;}
}
{
}

\Verse{11}
{
    \zA{二人正说着，只见香菱、臻儿、司棋、侍书等上亭子来了。 二人只得掩住这话，且和他们玩笑。 只见凤姐儿站在山坡上招手儿，小红便连忙弃
        了众人，跑至凤姐前，堆着笑问：“奶奶使唤做什么事？” 凤姐打量了一回，见他 生的干净俏丽，说话知趣，因笑道：“我的丫头们今儿没跟进我来。 我这会子想起
        一件事来，要使唤个人出去，不知你能干不能干？ 说的齐全不齐全？” 小红笑道： “奶奶有什么话，只管吩咐我说去；}
}
{
    \eA{and throughout the whole expanse of the park, embroidered sashes waved to
        and fro, and ornamented branches nodded their heads about. In addition to
        this, the members of the family were clad in such fineries that they put the
        peach tree to shame, made the almond yield the palm, the swallow envious and
        the hawk to blush. We could not therefore exhaustively describe them within
        our limited space of time.}
}
{
}

\Verse{12}
{
    \zA{要说的不齐全，误了奶奶的事，任凭奶奶责罚 就是了。” 凤姐笑道：“你是那位姑娘屋里的？ 我使你出去，他回来找你，我好替 你说。”
        小红道：“我是宝二爷屋里的。” 凤姐听了笑道：“嗳哟!你原来是宝玉 屋里的，怪道呢。 也罢了，等他问，我替你说。 你到我们家告诉你平姐姐，外头屋
        里桌子上汝窑盘子架儿底下放着一卷银子，那是一百二十两，给绣匠的工价。 等张 材家的来，当面秤给他瞧了，再给他拿去。}
}
{
    \eA{Pao-ch'ai, Ying Ch'un, T'an Ch'un, Hsi Ch'un, Li Wan, lady Feng and other
        girls, as well as Ta Chieh Erh, Hsiang Ling and the waiting-maids were, one
        and all, we will now notice, in the garden enjoying themselves; the only
        person who could not be seen was Lin Tai-yü. "How is it," consequently
        inquired Ying Ch'un, "that I don't see cousin Liu? What a lazy girl!}
}
{
}

\Verse{13}
{
    \zA{还有一件事：里头床头儿上有个小荷包 儿，拿了来。” 小红听说，答应着，撤身去了。 不多时回来，不见凤姐在山坡上了，因见司棋从山洞里出来，站着系带子，便
        赶来问道：“姐姐，不知道二奶奶往那里去了？” 司棋道：“没理论。” 小红听了， 回身又往四下里一看，只见那边探春宝钗在池边看鱼，小红上来陪笑道：“姑娘们
        可知道二奶奶刚才那里去了？” 探春道：“往你大奶奶院里找去。”}
}
{
    \eA{Is she forsooth fast asleep even at this late hour of the day?" "Wait all of
        you here," rejoined Pao-ch'ai, "and I'll go and shake her up and bring her."
        With these words, she speedily left her companions and repaired straightway
        into the Hsiao Hsiang lodge. While she was going on her errand, she met Wen
        Kuan and the rest of the girls, twelve in all, on their way to seek the
        party.}
}
{
}

\Verse{14}
{
    \zA{小红听了，再 往稻香村来，顶头见晴雯、绮霞、碧痕、秋纹、麝月、侍书、入画、莺儿等一群人 来了。
        晴雯一见小红，便说道：“你只是疯罢!院子里花儿也不浇，雀儿也不喂， 茶炉子也不弄，就在外头逛！” 小红道：“昨儿二爷说了，今儿不用浇花儿，过一
        日浇一回。 我喂雀儿的时候儿，你还睡觉呢。” 碧痕道：“茶炉子呢？” 小红道： “今儿不该我的班儿，有茶没茶，别问我。”
        绮霞道：“你听听他的嘴!你们别说 了，让他逛罢。”}
}
{
    \eA{Drawing near, they inquired after her health. After exchanging a few
        commonplace remarks, Pao-ch'ai turned round and pointing, said: "you will
        find them all in there; you had better go and join them. As for me, I'm
        going to fetch Miss Lin, but I'll be back soon." Saying this, she followed
        the winding path, and came to the Hsiao Hsiang lodge. Upon suddenly raising
        her eyes, she saw Pao-yü walk in.}
}
{
}

\Verse{15}
{
    \zA{小红道：“你们再问问，我逛了没逛。 二奶奶才使唤我说话取东 西去。” 说着，将荷包举给他们看，方没言语了，大家走开。 晴雯冷笑道：“怪道
        呢!原来爬上高枝儿去了，就不服我们说了。 不知说了一句话半句话，名儿姓儿知 道了没有，就把他兴头的这个样儿。 这一遭半遭儿的也算不得什么：过了后儿，还
        得听呵。 有本事从今儿出了这园子，长长远远的在高枝儿上才算好的呢！” 一面说 着去了。 这里小红听了，不便分证，只得忍气来找凤姐。}
}
{
    \eA{Pao-ch'ai immediately halted, and, lowering her head, she gave way to
        meditation for a time. "Pao-yü and Lin Tai-yü," she reflected, "have grown
        up together from their very infancy. But cousins, though they be, there are
        many instances in which they cannot evade suspicion, for they joke without
        heeding propriety; and at one time they are friends and at another at
        daggers drawn.}
}
{
}

\Verse{16}
{
    \zA{到了李氏房中，果见凤姐在这 里和李氏说话儿呢。 小红上来回道：“平姐姐说：奶奶刚出来了，他就把银子收起 来了； 才张材家的来取，当面秤了给他拿了去了。”
        说着，将荷包递上去。 又道： “平姐姐叫我来回奶奶：才旺儿进来讨奶奶的示下，好往那家子去，平姐姐就把那 话按着奶奶的主意打发他去了。”
        凤姐笑道：“他怎么按着我的主意打发去了呢？” 小红道：“平姐姐说：‘我们奶奶问这里奶奶好。 我们二爷没在家。 虽然迟了两天， 只管请奶奶放心。}
}
{
    \eA{Tai-yü has, moreover, always been full of envy; and has ever displayed a
        peevish disposition, so were I to follow him in at this juncture, why,
        Pao-yü would, in the first place, not feel at ease, and, in the second,
        Tai-yü would give way to jealousy. Better therefore for me to turn back." At
        the close of this train of thought, she retraced her steps.}
}
{
}

\Verse{17}
{
    \zA{等五奶奶好些，我们奶奶还会了五奶奶来瞧奶奶呢。 五奶奶前儿 打发了人来说：舅奶奶带了信来了，问奶奶好，还要和这里的姑奶奶寻几丸延年神 验万金丹；
        若有了，奶奶打发人来，只管送在我们奶奶这里。 明儿有人去，就顺路 给那边舅奶奶带了去。’” 小红还未说完，李氏笑道：“嗳哟!这话我就不懂了，
        什么‘奶奶’‘爷爷’的一大堆。” 凤姐笑道：“怨不得你不懂，这是四五门子的 话呢。”}
}
{
    \eA{But just as she was starting to join her other cousins, she unexpectedly
        descried, ahead of her, a pair of jade-coloured butterflies, of the size of
        a circular fan. Now they soared high, now they made a swoop down, in their
        flight against the breeze; much to her amusement.}
}
{
}

\Verse{18}
{
    \zA{说着，又向小红笑道：“好孩子，难为你说的齐全，不像他们扭扭捏捏蚊 子似的。 嫂子不知道，如今除了我随手使的这几个丫头老婆之外，我就怕和别人说
        话：他们必定把一句话拉长了，作两三截儿，咬文嚼字，拿着腔儿，哼哼唧唧的。 急的我冒火，他们那里知道？ 我们平儿先也是这么着，我就问着他：难道必定装蚊
        子哼哼就算美人儿了？ 说了几遭儿才好些儿了。” 李纨笑道：“都像你泼辣货才好。” 凤姐道：“这个丫头就好。 刚才这两遭说话虽不多，口角儿就很剪断。”}
}
{
    \eA{Pao-ch'ai felt a wish to catch them for mere fun's sake, so producing a fan
        from inside her sleeve, she descended on to the turfed ground to flap them
        with it. The two butterflies suddenly were seen to rise; suddenly to drop:
        sometimes to come; at others to go.}
}
{
}

\Verse{19}
{
    \zA{说着，又 向小红笑道：“明儿你伏侍我罢，我认你做干女孩儿。 我一调理，你就出息了。” 小红听了，“扑哧”一笑。 凤姐道：“你怎么笑？
        你说我年轻，比你能大几岁， 就做你的妈了？ 你做春梦呢!你打听打听，这些人比你大的赶着我叫妈，我还不理 呢，今儿抬举了你了。”
        小红笑道：“我不是笑这个，我笑奶奶认错了辈数儿了。 我妈是奶奶的干女孩儿，这会子又认我做干女孩儿！” 凤姐道：“谁是你妈？” 李
        纨笑道：“你原来不认的他？}
}
{
    \eA{Just as they were on the point of flying across the stream to the other
        side, the enticement proved too much for Pao-ch'ai, and she pursued them on
        tiptoe straight up to the Ti Ts'ui pavilion, nestling on the bank of the
        pond; while fragrant perspiration dripped drop by drop, and her sweet breath
        panted gently.}
}
{
}

\Verse{20}
{
    \zA{他是林之孝的女孩儿。” 凤姐听了，十分诧异，因说 道：“哦，是他的丫头啊。” 又笑道：“林子孝两口子，都是锥子扎不出一声儿来 的。
        我成日家说，他们倒是配就了的一对儿：一个‘天聋’，一个‘地哑’。 那里 承望养出这么个伶俐丫头来!你十几了？” 小红道：“十七岁了。” 又问名字。 小
        红道：“原叫‘红玉’，因为重了宝二爷，如今只叫小红了。” 凤姐听说，将眉一
        皱，把头一回，说道：“讨人嫌的很!得了‘玉’的便宜似的，你也‘玉’我也‘玉’。”}
}
{
    \eA{But Pao-ch'ai abandoned the idea of catching them, and was about to beat a
        retreat, when all at once she overheard, in the pavilion, the chatter of
        people engaged in conversation. This pavilion had, it must be added, a
        verandah and zig-zag balustrades running all round. It was erected over the
        water, in the centre of a pond, and had on the four sides window-frames of
        carved wood work, stuck with paper.}
}
{
}

\Verse{21}
{
    \zA{因说：“嫂子不知道，我和他妈说：‘赖大家的如今事多，也不知这府里谁是谁， 你替我好好儿的挑两个丫头我使。’ 他只管答应着； 他饶不挑，倒把他的女孩儿送
        给别处去。 难道跟我必定不好？” 李纨笑道：“你可是又多心了。 进来在先，你说 在后，怎么怨的他妈？”
        凤姐也笑道：“既这么着，明儿我和宝玉说，叫他再要人， 叫这丫头跟我去。 可不知本人愿意不愿意？” 小红笑道：“愿意不愿意，我们也不 敢说。}
}
{
    \eA{So when Pao-ch'ai caught, from without the pavilion, the sound of voices,
        she at once stood still and lent an attentive ear to what was being said.
        "Look at this handkerchief," she overheard. "If it's really the one you've
        lost, well then keep it; but if it isn't you must return it to Mr. Yün." "To
        be sure it is my own," another party observed, "bring it along and give it
        to me." "What reward will you give me?"}
}
{
}

\Verse{22}
{
    \zA{只是跟着奶奶，我们学些眉眼高低，出入上下，大小的事儿，也得见识见识。” 刚说着，只见王夫人的丫头来请，凤姐便辞了李纨去了。 小红自回怡红院去，不在
        话下。 如今且说黛玉因夜间失寝，次日起来迟了，闻得众姐妹都在园中做饯花会，恐 人笑他痴懒，连忙梳洗了出来。 刚到了院中，只见宝玉进门，来了便笑道：“好妹
        妹，你昨儿告了我了没有？ 叫我悬了一夜的心。”}
}
{
    \eA{she further heard. "Is it likely that I've searched all for nothing!" "I've
        long ago promised to recompense you, and of course I won't play you false,"
        some one again rejoined. "I found it and brought it round," also reached her
        ear, "and you naturally will recompense me; but won't you give anything to
        the person who picked it up?"}
}
{
}

\Verse{23}
{
    \zA{黛玉便回头叫紫鹃：“把屋子收 拾了，下一扇纱屉子，看那大燕子回来，把帘子放下来，拿狮子倚住。 烧了香，就 把炉罩上。” 一面说，一面又往外走。
        宝玉见他这样，还认作是昨日晌午的事，那 知晚间的这件公案？ 还打恭作揖的。 黛玉正眼儿也不看，各自出了院门，一直找别 的姐妹去了。
        宝玉心中纳闷，自己猜疑：“看起这样光景来，不像是为昨儿的事。 但只昨日我回来的晚了，又没有见他，再没有冲撞他的去处儿了。” 一面想，一面
        由不得随后跟了来。}
}
{
    \eA{"Don't talk nonsense," the other party added, "he belongs to a family of
        gentlemen, and anything of ours he may pick up it's his bounden duty to
        restore to us. What reward could you have me give him?" "If you don't reward
        him," she heard some one continue, "what will I be able to tell him?
        Besides, he enjoined me time after time that if there was to be no
        recompense, I was not to give it to you."}
}
{
}

\Verse{24}
{
    \zA{只见宝钗探春正在那边看鹤舞，见黛玉来了，三个一同站着说话儿。 又见宝玉 来了，探春便笑道：“宝哥哥身上好？ 我整整的三天没见你了。” 宝玉笑道：“妹
        妹身上好？ 我前儿还在大嫂子跟前问你呢。” 探春道：“宝哥哥，你往这里来，我 和你说话。” 宝玉听说，便跟了他，离了钗玉两个，到了一棵石榴树下。
        探春因说 道：“这几天，老爷没叫你吗？” 宝玉笑道：“没有叫。” 探春道：“昨儿我恍惚 听见说，老爷叫你出去来着。”}
}
{
    \eA{A short pause ensued. "Never mind!" then came out again to her, "take this
        thing of mine and present it to him and have done! But do you mean to let
        the cat out of the bag with any one else? You should take some oath." "If I
        tell any one," she likewise overheard, "may an ulcer grow on my mouth, and
        may I, in course of time, die an unnatural death!" "Ai-ya!" was the reply
        she heard;}
}
{
}

\Verse{25}
{
    \zA{宝玉笑道：“那想是别人听错了，并没叫我。” 探 春又笑道：“这几个月，我又攒下有十来吊钱了。 你还拿了去，明儿出门逛去的时
        候，或是好字画，好轻巧玩意儿，替我带些来。” 宝玉道：“我这么逛去，城里城 外大廊大庙的逛，也没见个新奇精致东西，总不过是那些金、玉、铜、磁器，没处
        撂的古董儿，再么就是绸缎、吃食、衣服了。”}
}
{
    \eA{"our minds are merely bent upon talking, but some one might come and quietly
        listen from outside; wouldn't it be as well to push all the venetians open.
        Any one seeing us in here will then imagine that we are simply chatting
        about nonsense. Besides, should they approach, we shall be able to observe
        them, and at once stop our conversation!"}
}
{
}

\Verse{26}
{
    \zA{探春道：“谁要那些作什么!像你 上回买的那柳枝儿编的小篮子儿，竹子根儿挖的香盒儿，胶泥垛的风炉子儿，就好 了，我喜欢的了不的。
        谁知他们都爱上了，都当宝贝儿似的抢了去了。” 宝玉笑道： “原来要这个。 这不值什么，拿几吊钱出去给小子们，管拉两车来。” 探春道：“小 厮们知道什么？
        你拣那有意思儿又不俗气的东西，你多替我带几件来，我还像上回 的鞋做一双你穿，比那双还加工夫，如何呢？”}
}
{
    \eA{Pao-ch'ai listened to these words from outside, with a heart full of
        astonishment. "How can one wonder," she argued mentally, "if all those lewd
        and dishonest people, who have lived from olden times to the present, have
        devised such thorough artifices! But were they now to open and see me here,
        won't they feel ashamed.}
}
{
}

\Verse{27}
{
    \zA{宝玉笑道：“你提起鞋来，我想起故事来了：一回穿着，可巧遇见了老爷，老 爷就不受用，问：‘是谁做的？’ 我那里敢提三妹妹，我就回说是前儿我的生日舅
        母给的。 老爷听了是舅母给的，才不好说什么了。 半日还说：‘何苦来!虚耗人力， 作践绫罗，做这样的东西。’
        我回来告诉了袭人，袭人说：‘这还罢了，赵姨娘气 的抱怨的了不得：正经亲兄弟，鞋塌拉袜塌拉的没人看见，且做这些东西！’”}
}
{
    \eA{Moreover, the voice in which those remarks were uttered resembles very much
        that of Hung Erh, attached to Pao-yü's rooms, who has all along shown a
        sharp eye and a shrewd mind. She's an artful and perverse thing of the first
        class!}
}
{
}

\Verse{28}
{
    \zA{探 春听说，登时沉下脸来，道：“你说，这话糊涂到什么田地!怎么我是该做鞋的人 么？ 环儿难道没有分例的？
        衣裳是衣裳，鞋袜是鞋袜，丫头老婆一屋子，怎么抱怨这 些话？ 给谁听呢!我不过闲着没事作一双半双，爱给那个哥哥兄弟，随我的心，谁敢 管我不成？
        这也是他瞎气。” 宝玉听了，点头笑道：“你不知道，他心里自然又有 个想头了。” 探春听说，一发动了气，将头一扭，说道：“连你也糊涂了!他那想头，自然
        是有的。 不过是那阴微下贱的见识。}
}
{
    \eA{And as I have now overheard her peccadilloes, and a person in despair rebels
        as sure as a dog in distress jumps over the wall, not only will trouble
        arise, but I too shall derive no benefit. It would be better at present
        therefore for me to lose no time in retiring.}
}
{
}

\Verse{29}
{
    \zA{他只管这么想，我只管认得老爷太太两个人， 别人我一概不管。 就是姐妹弟兄跟前，谁和我好，我就和谁好； 什么偏的庶的，我 也不知道。
        论理我不该说他，但他忒昏愦的不像了!还有笑话儿呢：就是上回我给 你那钱，替我买那些玩的东西，过了两天，他见了我，就说是怎么没钱，怎么难过。 我也不理。
        谁知后来丫头们出去了，他就抱怨起我来，说我攒的钱为什么给你使， 倒不给环儿使呢!我听见这话，又好笑又好气。 我就出来往太太跟前去了。”}
}
{
    \eA{But as I fear I mayn't be in time to get out of the way, the only
        alternative for me is to make use of some art like that of the cicada, which
        can divest itself of its \_exuviae\_." She had scarcely brought her
        reflections to a close before a sound of 'ko-chih' reached her ears.
        Pao-ch'ai purposely hastened to tread with heavy step. "P'in Erh, I see
        where you're hiding!" she cried out laughingly;}
}
{
}

\Verse{30}
{
    \zA{正说 着，只见宝钗那边笑道：“说完了？ 来罢。 显见的是哥哥妹妹了，撂下别人，且说 体己去。 我们听一句儿就使不得了？” 说着，探春宝玉二人方笑着来了。
        宝玉因不见了黛玉，便知是他躲了别处去了。 想了一想：“索性迟两日，等他 的气息一息再去也罢了。” 因低头看见许多凤仙石榴等各色落花，锦重重的落了一
        地，因叹道：“这是他心里生了气，也不收拾这花儿来了。 等我送了去，明儿再问 着他。” 说着，只见宝钗约着他们往后头去。}
}
{
    \eA{and as she shouted, she pretended to be running ahead in pursuit of her. As
        soon as Hsiao Hung and Chui Erh pushed the windows open from inside the
        pavilion, they heard Pao-ch'ai screaming, while rushing forward; and both
        fell into a state of trepidation from the fright they sustained. Pao-ch'ai
        turned round and faced them. "Where have you been hiding Miss Lin?" she
        smiled.}
}
{
}

\Verse{31}
{
    \zA{宝玉道：“我就来。” 等他二人去远， 把那花儿兜起来，登山渡水，过树穿花，一直奔了那日和黛玉葬桃花的去处。
        将已到了花冢，犹未转过山坡，只听那边有呜咽之声，一面数落着，哭的好不 伤心。 宝玉心下想道：“这不知是那屋里的丫头，受了委屈，跑到这个地方来哭？”
        一面想，一面煞住脚步，听他哭道是： 花谢花飞飞满天，红消香断有谁怜？ 游丝软系飘春榭，落絮轻沾扑绣帘。 闺中女儿惜春暮，愁绪满怀无着处。
        手把花锄出绣帘，忍踏落花来复去？}
}
{
    \eA{"Who has seen anything of Miss Lin," retorted Chui Erh. "I was just now,"
        proceeded Pao-ch'ai, "on that side of the pool, and discerned Miss Lin
        squatting down over there and playing with the water. I meant to have gently
        given her a start, but scarcely had I walked up to her, when she saw me,
        and, with a \_detour\_ towards the East, she at once vanished from sight. So
        mayn't she be concealing herself in there?"}
}
{
}

\Verse{32}
{
    \zA{柳丝榆荚自芳菲，不管桃飘与李飞。 桃李明年能再发，明年闺中知有谁？ 三月香巢初垒成，梁间燕子太无情! 明年花发虽可啄，却不道人去梁空巢已倾。
        一年三百六十日，风刀霜剑严相逼。 明媚鲜妍能几时，一朝飘泊难寻觅。 花开易见落难寻，阶前愁杀葬花人。 独把花锄偷洒泪，洒上空枝见血痕。
        杜鹃无语正黄昏，荷锄归去掩重门。 青灯照壁人初睡，冷雨敲窗被未温。 怪侬底事倍伤神？ 半为怜春半恼春： 怜春忽至恼忽去，至又无言去不闻。}
}
{
    \eA{As she spoke, she designedly stepped in and searched about for her. This
        over, she betook herself away, adding: "she's certain to have got again into
        that cave in the hill, and come across a snake, which must have bitten her
        and put an end to her." So saying, she distanced them, feeling again very
        much amused.}
}
{
}

\Verse{33}
{
    \zA{昨宵庭外悲歌发，知是花魂与鸟魂？ 花魂鸟魂总难留，鸟自无言花自羞。 愿侬此日生双翼，随花飞到天尽头。 天尽头，何处有香丘？ 未若锦囊收艳骨，一？
        净土掩风流。 质本洁来还洁去，不教污淖陷渠沟。 尔今死去侬收葬，未卜侬身何日丧？ 侬今葬花人笑痴，他年葬侬知是谁？ 试看春残花渐落，便是红颜老死时。
        一朝春尽红颜老，花落人亡两不知! 正是一面低吟，一面哽咽。 那边哭的自己伤心，却不道这边听的早已疾倒了。 要知端详，下回分解。 (本章完)}
}
{
    \eA{"I have managed," she thought, "to ward off this piece of business, but I
        wonder what those two think about it." Hsiao Hung, who would have
        anticipated, readily credited as gospel the remarks she heard Pao-ch'ai
        make. But allowing just time enough to Pao-ch'ai to got to a certain
        distance, she instantly drew Chui Erh to her. "Dreadful!"}
}
{
}

\Verse{34}
{
    \zA{}
}
{
    \eA{she observed, "Miss Lin was squatting in here and must for a certainty have
        overheard what we said before she left." Albeit Chui Erh listened to her
        words, she kept her own counsel for a long time. "What's to be done?" Hsiao
        Hung consequently exclaimed. "Even supposing she did overhear what we said,"
        rejoined Chui Erh by way of answer, "why should she meddle in what does not
        concern her? Every one should mind her own business." "Had it been Miss Pao,
        it would not have mattered," remarked Hsiao Hung, "but Miss Lin delights in
        telling mean things of people and is, besides, so petty-minded. Should she
        have heard and anything perchance comes to light, what will we do?"}
}
{
}

\Verse{35}
{
    \zA{}
}
{
    \eA{During their colloquy, they noticed Wen Kuan, Hsiang Ling, Ssu Ch'i, Shih
        Shu and the other girls enter the pavilion, so they were compelled to drop
        the conversation and to play and laugh with them. They then espied lady Feng
        standing on the top of the hillock, waving her hand, beckoning to Hsiao
        Hung. Hurriedly therefore leaving the company, she ran up to lady Feng and
        with smile heaped upon smile, "my lady," she inquired, "what is it that you
        want?" Lady Feng scrutinised her for a time. Observing how spruce and pretty
        she was in looks, and how genial in her speech, she felt prompted to give
        her a smile. "My own waiting-maid," she said, "hasn't followed me in here
        to-day;}
}
{
}

\Verse{36}
{
    \zA{}
}
{
    \eA{and as I've just this moment bethought myself of something and would like to
        send some one on an errand, I wonder whether you're fit to undertake the
        charge and deliver a message faithfully." "Don't hesitate in entrusting me
        with any message you may have to send," replied Hsiao Hung with a laugh.
        "I'll readily go and deliver it. Should I not do so faithfully, and blunder
        in fulfilling your business, my lady, you may visit me with any punishment
        your ladyship may please, and I'll have nothing to say." "What young lady's
        servant are you," smiled lady Feng? "Tell me, so that when she comes back,
        after I've sent you out, and looks for you, I may be able to tell her about
        you."}
}
{
}

\Verse{37}
{
    \zA{}
}
{
    \eA{"I'm attached to our Master Secundus,' Mr. Pao's rooms," answered Hsiao
        Hung. "Ai-ya!" ejaculated lady Feng, as soon as she heard these words. "Are
        you really in Pao-yü's rooms! How strange! Yet it comes to the same thing.
        Well, if he asks for you, I'll tell him where you are. Go now to our house
        and tell your sister P'ing that she'll find on the table in the outer
        apartment and under the stand with the plate from the Ju kiln, a bundle of
        silver; that it contains the one hundred and twenty taels for the
        embroiderers' wages; and that when Chang Ts'ai's wife comes, the money
        should be handed to her to take away, after having been weighed in her
        presence and been given to her to tally.}
}
{
}

\Verse{38}
{
    \zA{}
}
{
    \eA{Another thing too I want. In the inner apartment and at the head of the bed
        you'll find a small purse, bring it along to me." Hsiao Hung listened to her
        orders and then started to carry them out. On her return, in a short while,
        she discovered that lady Feng was not on the hillock. But perceiving Ssu
        Ch'i egress from the cave and stand still to tie her petticoat, she walked
        up to her. "Sister, do you know where our lady Secunda is gone to?" she
        asked. "I didn't notice," rejoined Ssu Ch'i. At this reply, Hsiao Hung
        turned round and cast a glance on all four quarters. Seeing T'an Ch'un and
        Pao-ch'ai standing by the bank of the pond on the opposite side and looking
        at the fish, Hsiao Hung advanced up to them.}
}
{
}

\Verse{39}
{
    \zA{}
}
{
    \eA{"Young ladies," she said, straining a smile, "do you perchance have any idea
        where our lady Secunda is gone to now?" "Go into your senior lady's court
        and look for her!" T'an Ch'un answered. Hearing this, Hsiao Hung was
        proceeding immediately towards the Tao Hsiang village, when she caught
        sight, just ahead of her, of Ch'ing Wen, Ch'i Hsia, Pi Hen, Ch'iu Wen, She
        Yüeh, Shih Shu, Ju Hua, Ying Erh and some other girls coming towards her in
        a group. The moment Ch'ing Wen saw Hsiao Hung, she called out to her. "Are
        you gone clean off your head?" she exclaimed.}
}
{
}

\Verse{40}
{
    \zA{}
}
{
    \eA{"You don't water the flowers, nor feed the birds or prepare the tea stove,
        but gad about outside!" "Yesterday," replied Hsiao Hung, "Mr. Secundus told
        me that there was no need for me to water the flowers to-day; that it was
        enough if they were watered every other day. As for the birds, you're still
        in the arms of Morpheus, sister, when I give them their food." "And what
        about the tea-stove?" interposed Pi Hen. "To-day," retorted Hsiao Hung, "is
        not my turn on duty, so don't ask me whether there be any tea or not!" "Do
        you listen to that mouth of hers!" cried Ch'i Hsia, "but don't you girls
        speak to her; let her stroll about and have done!" "You'd better all go and
        ask whether I've been gadding about or not," continued Hsiao Hung. "Our lady
        Secunda has just bidden me go and deliver a message, and fetch something."}
}
{
}

\Verse{41}
{
    \zA{}
}
{
    \eA{Saying this, she raised the purse and let them see it; and they, finding
        they could hit upon nothing more to taunt her with, trudged along onwards.
        Ch'ing Wen smiled a sarcastic smile. "How funny!" she cried. "Lo, she climbs
        up a high branch and doesn't condescend to look at any one of us! All she
        told her must have been just some word or two, who knows! But is it likely
        that our lady has the least notion of her name or surname that she rides
        such a high horse, and behaves in this manner! What credit is it in having
        been sent on a trifling errand like this! Will we, by and bye, pray, hear
        anything more about you? If you've got any gumption, you'd better skedaddle
        out of this garden this very day.}
}
{
}

\Verse{42}
{
    \zA{}
}
{
    \eA{For, mind, it's only if you manage to hold your lofty perch for any length
        of time that you can be thought something of!" As she derided her, she
        continued on her way. During this while, Hsiao Hung listened to her, but as
        she did not find it a suitable moment to retaliate, she felt constrained to
        suppress her resentment and go in search of lady Feng. On her arrival at
        widow Li's quarters, she, in point of fact, discovered lady Feng seated
        inside with her having a chat. Hsiao Hung approached her and made her
        report.}
}
{
}

\Verse{43}
{
    \zA{}
}
{
    \eA{"Sister P'ing says," she observed, "that as soon as your ladyship left the
        house, she put the money by, and that when Chang Ts'ai's wife went in a
        little time to fetch it, she had it weighed in her presence, after which she
        gave it to her to take away." With these words, she produced the purse and
        presented it to her. "Sister P'ing bade me come and tell your ladyship," she
        added, continuing, "that Wang Erh came just now to crave your orders, as to
        who are the parties from whom he has to go and (collect interest on money
        due) and sister P'ing explained to him what your wishes were and sent him
        off." "How could she tell him where I wanted him to go?" Lady Feng laughed.}
}
{
}

\Verse{44}
{
    \zA{}
}
{
    \eA{"Sister P'ing says," Hsiao Hung proceeded, "that our lady presents her
        compliments to your ladyship (widow Li) here-(\_To lady Feng\_) that our
        master Secundus has in fact not come home, and that albeit a delay of (a
        day) or two will take place (in the collection of the money), your ladyship
        should, she begs, set your mind at ease. (\_To Li Wan\_). That when lady
        Quinta is somewhat better, our lady will let lady Quinta know and come along
        with her to see your ladyship. (\_To lady Feng\_). That lady Quinta sent a
        servant the day before yesterday to come over and say that our lady, your
        worthy maternal aunt, had despatched a letter to inquire after your
        ladyship's health;}
}
{
}

\Verse{45}
{
    \zA{}
}
{
    \eA{that she also wished to ask you, my lady, her worthy niece in here, for a
        couple of 'long-life-great-efficacy-full-of-every-virtue' pills; and that if
        you have any, they should, when our lady bids a servant come over, be simply
        given her to bring to our lady here, and that any one bound to-morrow for
        that side could then deliver them on her way to her ladyship, your aunt
        yonder, to take along with her." "Ai-yo-yo!" exclaimed widow Li, before the
        close of the message. "It's impossible for me to make out what you're
        driving at! What a heap of ladyships and misters!" "It's not to be wondered
        at that you can't make them out," interposed lady Feng laughing. "Why, her
        remarks refer to four or five distinct families." While speaking, she again
        faced Hsiao Hung. "My dear girl," she smiled, "what a trouble you've been
        put to!}
}
{
}

\Verse{46}
{
    \zA{}
}
{
    \eA{But you speak decently, and unlike the others who keep on buzz-buzz-buzz,
        like mosquitoes! You're not aware, sister-in-law, that I actually dread
        uttering a word to any of the girls outside the few servant-girls and
        matrons in my own immediate service; for they invariably spin out, what
        could be condensed in a single phrase, into a long interminable yarn, and
        they munch and chew their words; and sticking to a peculiar drawl, they
        groan and moan; so much so, that they exasperate me till I fly into a
        regular rage. Yet how are they to know that our P'ing Erh too was once like
        them. But when I asked her: 'must you forsooth imitate the humming of a
        mosquito, in order to be accounted a handsome girl?' and spoke to her, on
        several occasions, she at length improved considerably."}
}
{
}

\Verse{47}
{
    \zA{}
}
{
    \eA{"What a good thing it would be," laughed Li Kung-ts'ai, "if they could all
        be as smart as you are." "This girl is first-rate!" rejoined lady Feng, "she
        just now delivered two messages. They didn't, I admit, amount to much, yet
        to listen to her, she spoke to the point." "To-morrow," she continued,
        addressing herself to Hsiao Hung smilingly, "come and wait on me, and I'll
        acknowledge you as my daughter; and the moment you come under my control,
        you'll readily improve." At this news, Hsiao Hung spurted out laughing
        aloud. "What are you laughing for?" Lady Feng inquired. "You must say to
        yourself that I am young in years and that how much older can I be than
        yourself to become your mother; but are you under the influence of a spring
        dream? Go and ask all those people older than yourself.}
}
{
}

\Verse{48}
{
    \zA{}
}
{
    \eA{They would be only too ready to call me mother. But snapping my fingers at
        them, I to-day exalt you." "I wasn't laughing about that," Hsiao Hung
        answered with a smiling face. "I was amused by the mistake your ladyship
        made about our generations. Why, my mother claims to be your daughter, my
        lady, and are you now going to recognise me too as your daughter?" "Who's
        your mother?" Lady Feng exclaimed. "Don't you actually know her?" put in Li
        Kung-ts'ai with a smile. "She's Lin Chih-hsiao's child." This disclosure
        greatly surprised lady Feng. "What!" she consequently cried, "is she really
        his daughter?" "Why Lin Chih-hsiao and his wife," she resumed smilingly,
        "couldn't either of them utter a sound if even they were pricked with an
        awl. I've always maintained that they're a well-suited couple;}
}
{
}

\Verse{49}
{
    \zA{}
}
{
    \eA{as the one is as deaf as a post, and the other as dumb as a mute. But who
        would ever have expected them to have such a clever girl! By how much are
        you in your teens?" "I'm seventeen," replied Hsia Hung. "What is your name?"
        she went on to ask. "My name was once Hung Yü." Hsiao Hung rejoined. "But as
        it was a duplicate of that of Master Secundus, Mr. Pao-yü, I'm now simply
        called Hsiao Hung." Upon hearing this explanation, lady Feng raised her
        eyebrows into a frown, and turning her head round: "It's most disgusting!"
        she remarked, "Those bearing the name Yü would seem to be very cheap; for
        your name is Yü, and so is also mine Yü. Sister-in-law," she then observed;}
}
{
}

\Verse{50}
{
    \zA{}
}
{
    \eA{"I never let you know anything about it, but I mentioned to her mother that
        Lai Ta's wife has at present her hands quite full, and that she hasn't
        either any notion as to who is who in this mansion. 'You had better,' (I
        said), 'carefully select a couple of girls for my service.' She assented
        unreservedly, but she put it off and never chose any. On the contrary, she
        sent this girl to some other place. But is it likely that she wouldn't have
        been well off with me?" "Here you are again full of suspicion!" Li Wan
        laughed. "She came in here long before you ever breathed a word to her! So
        how could you bear a grudge against her mother?" "Well, in that case," added
        lady Feng, "I'll speak to Pao-yü to-morrow, and induce him to find another
        one, and to allow this girl to come along with me.}
}
{
}

\Verse{51}
{
    \zA{}
}
{
    \eA{I wonder, however, whether she herself is willing or not?" "Whether willing
        or not," interposed Hsiao Hung smiling, "such as we couldn't really presume
        to raise our voices and object. We should feel it our privilege to serve
        such a one as your ladyship, and learn a little how to discriminate when
        people raise or drop their eyebrows and eyes (with pleasure or displeasure),
        and reap as well some experience in such matters as go out or come in,
        whether high or low, great and small." But during her reply, she perceived
        Madame Wang's waiting-maid come and invite lady Feng to go over.}
}
{
}

\Verse{52}
{
    \zA{}
}
{
    \eA{Lady Feng bade good-bye at once to Li Kung-ts'ai and took her departure.
        Hsiao Hung then returned into the I Hung court, where we will leave her and
        devote our attention for the present to Lin Tai-yü. As she had had but
        little sleep in the night, she got up the next day at a late hour. When she
        heard that all her cousins were collected in the park, giving a farewell
        entertainment for the god of flowers, she hastened, for fear people should
        laugh at her for being lazy, to comb her hair, perform her ablutions, and go
        out and join them. As soon as she reached the interior of the court, she
        caught sight of Pao-yü, entering the door, who speedily greeted her with a
        smile.}
}
{
}

\Verse{53}
{
    \zA{}
}
{
    \eA{"My dear cousin," he said, "did you lodge a complaint against me yesterday?
        I've been on pins and needles the whole night long." Tai-yü forthwith turned
        her head away. "Put the room in order," she shouted to Tzu Chüan, "and lower
        one of the gauze window-frames. And when you've seen the swallows come back,
        drop the curtain; keep it down then by placing the lion on it, and after you
        have burnt the incense, mind you cover the censer." So saying she stepped
        outside. Pao-yü perceiving her manner, concluded again that it must be on
        account of the incident of the previous noon, but how could he have had any
        idea about what had happened in the evening? He kept on still bowing and
        curtseying;}
}
{
}

\Verse{54}
{
    \zA{}
}
{
    \eA{but Lin Tai-yü did not even so much as look at him straight in the face, but
        egressing alone out of the door of the court, she proceeded there and then
        in search of the other girls. Pao-yü fell into a despondent mood and gave
        way to conjectures. "Judging," he reflected, "from this behaviour of hers,
        it would seem as if it could not be for what transpired yesterday. Yesterday
        too I came back late in the evening, and, what's more, I didn't see her, so
        that there was no occasion on which I could have given her offence." As he
        indulged in these reflections, he involuntarily followed in her footsteps to
        try and catch her up, when he descried Pao-ch'ai and T'an-ch'un on the
        opposite side watching the frolics of the storks.}
}
{
}

\Verse{55}
{
    \zA{}
}
{
    \eA{As soon as they saw Tai-yü approach, the trio stood together and started a
        friendly chat. But noticing Pao-yü also come up, T'an Ch'un smiled. "Brother
        Pao," she said, "are you all right. It's just three days that I haven't seen
        anything of you?" "Are you sister quite well?" Pao-yü rejoined, a smile on
        his lips. "The other day, I asked news of you of our senior sister-in-law."
        "Brother Pao," T'an Ch'un remarked, "come over here; I want to tell you
        something." The moment Pao-yü heard this, he quickly went with her.
        Distancing Pao-ch'ai and Tai-yü, the two of them came under a pomegranate
        tree. "Has father sent for you these last few days?" T'an Ch'un then asked.
        "He hasn't," Pao-yü answered laughingly by way of reply.}
}
{
}

\Verse{56}
{
    \zA{}
}
{
    \eA{"Yesterday," proceeded T'an Ch'un, "I heard vaguely something or other about
        father sending for you to go out." "I presume," Pao-yü smiled, "that some
        one must have heard wrong, for he never sent for me." "I've again managed to
        save during the last few months," added T'an Ch'un with another smile,
        "fully ten tiaos, so take them and bring me, when at any time you stroll out
        of doors, either some fine writings or some ingenious knicknack." "Much as I
        have roamed inside and outside the city walls," answered Pao-yü, "and seen
        grand establishments and large temples, I've never come across anything
        novel or pretty.}
}
{
}

\Verse{57}
{
    \zA{}
}
{
    \eA{One simply sees articles made of gold, jade, copper and porcelain, as well
        as such curios for which we could find no place here. Besides these, there
        are satins, eatables, and wearing apparel." "Who cares for such baubles!"
        exclaimed T'an Ch'un. "How could they come up to what you purchased the last
        time; that wee basket, made of willow twigs, that scent-box, scooped out of
        a root of real bamboo, that portable stove fashioned of glutinous clay;
        these things were, oh, so very nice! I was as fond of them as I don't know
        what; but, who'd have thought it, they fell in love with them and bundled
        them all off, just as if they were precious things." "Is it things of this
        kind that you really want?" laughed Pao-yü.}
}
{
}

\Verse{58}
{
    \zA{}
}
{
    \eA{"Why, these are worth nothing! Were you to take a hundred cash and give them
        to the servant-boys, they could, I'm sure, bring two cart-loads of them."
        "What do the servant-boys know?" T'an Ch'un replied. "Those you chose for me
        were plain yet not commonplace. Neither were they of coarse make. So were
        you to procure me as many as you can get of them, I'll work you a pair of
        slippers like those I gave you last time, and spend twice as much trouble
        over them as I did over that pair you have. Now, what do you say to this
        bargain?" "Your reference to this," smiled Pao-yü, "reminds me of an old
        incident.}
}
{
}

\Verse{59}
{
    \zA{}
}
{
    \eA{One day I had them on, and by a strange coincidence, I met father, whose
        fancy they did not take, and he inquired who had worked them. But how could
        I muster up courage to allude to the three words: my sister Tertia, so I
        answered that my maternal aunt had given them to me on the recent occasion
        of my birthday. When father heard that they had been given to me by my aunt,
        he could not very well say anything. But after a while, 'why uselessly
        waste,' he observed, 'human labour, and throw away silks to make things of
        this sort!' On my return, I told Hsi Jen about it. 'Never mind,' said Hsi
        Jen; but Mrs. Chao got angry. 'Her own brother,' she murmured indignantly,
        'wears slipshod shoes and socks in holes, and there's no one to look after
        him, and does she go and work all these things!'"}
}
{
}

\Verse{60}
{
    \zA{}
}
{
    \eA{T'an Ch'un, hearing this, immediately lowered her face. "Now tell me, aren't
        these words utter rot!" she shouted. "What am I that I have to make shoes?
        And is it likely that Huan Erh hasn't his own share of things! Clothes are
        clothes, and shoes and socks are shoes and socks; and how is it that any
        grudges arise in the room of a mere servant-girl and old matron? For whose
        benefit does she come out with all these things! I simply work a pair or
        part of a pair when I am at leisure, with time on my hands. And I can give
        them to any brother, elder or younger, I fancy; and who has a right to
        interfere with me? This is just another bit of blind anger!" After listening
        to her, Pao-yü nodded his head and smiled. "Yet," he said, "you don't know
        what her motives may be.}
}
{
}

\Verse{61}
{
    \zA{}
}
{
    \eA{It's but natural that she should also cherish some expectations." This
        apology incensed T'an Ch'un more than ever, and twisting her head round,
        "Even you have grown dull!" she cried. "She does, of course, indulge in
        expectations, but they are actuated by some underhand and paltry notion! She
        may go on giving way to these ideas, but I, for my part, will only care for
        Mr. Chia Cheng and Madame Wang. I won't care a rap for any one else. In
        fact, I'll be nice with such of my sisters and brothers, as are nice to me;
        and won't even draw any distinction between those born of primary wives and
        those of secondary ones. Properly speaking, I shouldn't say these things
        about her, but she's narrow-minded to a degree, and unlike what she should
        be. There's besides another ridiculous thing.}
}
{
}

\Verse{62}
{
    \zA{}
}
{
    \eA{This took place the last time I gave you the money to get me those trifles.
        Well, two days after that, she saw me, and she began again to represent that
        she had no money and that she was hard up. Nevertheless, I did not worry my
        brain with her goings on. But as it happened, the servant-girls subsequently
        quitted the room, and she at once started finding fault with me. 'Why,' she
        asked, 'do I give you my savings to spend and don't, after all, let Huan Erh
        have them and enjoy them?' When I heard these reproaches, I felt both
        inclined to laugh, and also disposed to lose my temper; but I there and then
        skedaddled out of her quarters, and went over to our Madame Wang."}
}
{
}

\Verse{63}
{
    \zA{}
}
{
    \eA{As she was recounting this incident, "Well," she overheard Pao-ch'ai
        sarcastically observe from the opposite direction, "have you done spinning
        your yarns? If you have, come along! It's quite evident that you are brother
        and sister, for here you leave every one else and go and discuss your own
        private matters. Couldn't we too listen to a single sentence of what you
        have to say?" While she taunted them, T'an Ch'un and Pao-yü eventually drew
        near her with smiling faces. Pao-yü, however, failed to see Lin Tai-yü and
        he concluded that she had dodged out of the way and gone elsewhere. "It
        would be better," he muttered, after some thought, "that I should let two
        days elapse, and give her temper time to evaporate before I go to her."}
}
{
}

\Verse{64}
{
    \zA{}
}
{
    \eA{But as he drooped his head, his eye was attracted by a heap of
        touch-me-nots, pomegranate blossom and various kinds of fallen flowers,
        which covered the ground thick as tapestry, and he heaved a sigh. "It's
        because," he pondered, "she's angry that she did not remove these flowers;
        but I'll take them over to the place, and by and bye ask her about them." As
        he argued to himself, he heard Pao-ch'ai bid them go out. "I'll join you in
        a moment," Pao-yü replied; and waiting till his two cousins had gone some
        distance, he bundled the flowers into his coat, and ascending the hill, he
        crossed the stream, penetrated into the arbour, passed through the avenues
        with flowers and wended his way straight for the spot, where he had, on a
        previous occasion, interred the peach-blossoms with the assistance of Lin
        Tai-yü.}
}
{
}

\Verse{65}
{
    \zA{}
}
{
    \eA{But scarcely had he reached the mound containing the flowers, and before he
        had, as yet, rounded the brow of the hill, than he caught, emanating from
        the off side, the sound of some one sobbing, who while giving way to
        invective, wept in a most heart-rending way. "I wonder," soliloquised
        Pao-yü, "whose servant-girl this is, who has been so aggrieved as to run
        over here to have a good cry!" While speculating within himself, he halted.
        He then heard, mingled with wails:—  Flowers wither and decay; and flowers
        do fleet; they fly all o'er the  skies; Their bloom wanes; their smell dies;
        but who is there with them to  sympathise?}
}
{
}

\Verse{66}
{
    \zA{}
}
{
    \eA{While vagrant gossamer soft doth on fluttering spring-bowers bind its
        coils,  And drooping catkins lightly strike and cling on the embroidered
        screens,  A maiden in the inner rooms, I sore deplore the close of spring.
        Such ceaseless sorrow fills my breast, that solace nowhere can I  find. Past
        the embroidered screen I issue forth, taking with me a hoe,  And on the
        faded flowers to tread I needs must, as I come and go. The willow fibres and
        elm seeds have each a fragrance of their own. What care I, peach blossoms
        may fall, pear flowers away be blown; Yet peach and pear will, when next
        year returns, burst out again in  bloom,  But can it e'er be told who will
        next year dwell in the inner room?}
}
{
}

\Verse{67}
{
    \zA{}
}
{
    \eA{What time the third moon comes, the scented nests have been already  built.
        And on the beams the swallows perch, excessive spiritless and staid; Next
        year, when the flowers bud, they may, it's true, have ample to  feed on:
        But they know not that when I'm gone beams will be vacant and nests  fall!
        In a whole year, which doth consist of three hundred and sixty days,  Winds
        sharp as swords and frost like unto spears each other rigorous  press,  So
        that how long can last their beauty bright; their fresh charm how  long
        stays? Sudden they droop and fly; and whither they have flown, 'tis hard to
        guess. Flowers, while in bloom, easy the eye attract; but, when they wither,
        hard they are to find.}
}
{
}

\Verse{68}
{
    \zA{}
}
{
    \eA{Now by the footsteps, I bury the flowers, but sorrow will slay me. Alone I
        stand, and as I clutch the hoe, silent tears trickle down,  And drip on the
        bare twigs, leaving behind them the traces of blood. The goatsucker hath
        sung his song, the shades lower of eventide,  So with the lotus hoe I return
        home and shut the double doors. Upon the wall the green lamp sheds its rays
        just as I go to sleep. The cover is yet cold; against the window patters the
        bleak rain. How strange! Why can it ever be that I feel so wounded at heart!
        Partly, because spring I regret; partly, because with spring I'm  vexed!
        Regret for spring, because it sudden comes; vexed, for it sudden  goes. For
        without warning, lo! it comes; and without asking it doth fleet.}
}
{
}

\Verse{69}
{
    \zA{}
}
{
    \eA{Yesterday night, outside the hall sorrowful songs burst from my  mouth,  For
        I found out that flowers decay, and that birds also pass away. The soul of
        flowers, and the spirit of birds are both hard to  restrain. Birds, to
        themselves when left, in silence plunge; and flowers,  alone,  they blush.
        Oh! would that on my sides a pair of wings could grow,  That to the end of
        heaven I may fly in the wake of flowers! Yea to the very end of heaven,
        Where I could find a fragrant grave! For better, is it not, that an
        embroidered bag should hold my  well-shaped bones,  And that a heap of
        stainless earth should in its folds my winsome  charms enshroud. For
        spotless once my frame did come, and spotless again it will go! Far better
        than that I, like filthy mire, should sink into some  drain!}
}
{
}

\Verse{70}
{
    \zA{}
}
{
    \eA{Ye flowers are now faded and gone, and, lo, I come to bury you. But as for
        me, what day I shall see death is not as yet divined! Here I am fain these
        flowers to inter; but humankind will laugh me as  a fool. Who knows, who
        will, in years to come, commit me to my grave! Mark, and you'll find the
        close of spring, and the gradual decay of  flowers,  Resemble faithfully the
        time of death of maidens ripe in years! In a twinkle, spring time draws to a
        close, and maidens wax in age. Flowers fade and maidens die; and of either
        nought any more is known. After listening to these effusions, Pao-yü
        unconsciously threw himself down in a wandering frame of mind. But, reader,
        do you feel any interest in him? If you do, the subsequent chapter contains
        further details about him.}
}
{
}
